\section{Conclusions\label{sec:conclusions}}

This study derived activity-based experience profiles for 71,467 GitHub developers linked to the AIDev dataset and examined how those profiles relate to AI-coding-agent use intensity and pull-request acceptance, alongside a separate comparison of agent-authored and human-authored pull-request acceptance. None of the three research questions shows a simple, monotonic relationship with activity level. Agent-use intensity is strongly non-monotonic: the intermediate-activity Mid profile submits far more Agentic-PRs than either Junior or the highest-activity Senior profile, and this is the one outcome with a large effect size (RQ1); acceptance rate is highest for Mid and lowest for Senior, with Junior in between (RQ2); and Agentic-PRs are accepted at a modestly higher rate than an independent human-authored baseline from a comparably popular subset of repositories (RQ3). Taken together, higher composite GitHub activity is not a reliable proxy for better or more frequent outcomes when using a coding agent: the highest-activity Senior profile has neither the most Agentic-PRs nor the highest acceptance rate of the three.

A second, more methodological conclusion follows from how this analysis was arrived at. Several of the reported associations changed direction, not just magnitude, at two separate points in the analysis: once when the underlying data sources were corrected (RQ2 initially appeared to favor Mid-activity developers because acceptance was computed on the wrong table), and again when the clustering preprocessing itself was corrected (a quantile-transformation step was found, via a complete-case sensitivity check, to be unstable to exactly which developers had complete GraphQL data -- Adjusted Rand Index $0.46$ against a $\log(1+x)$-based alternative's $0.98$ -- which changed RQ1 from a small, two-group effect into the largest effect size in the paper). Neither which table a query reads from, nor whether a per-developer indicator has full-population coverage, nor whether a transformation is stable to the exact composition of the sample, are incidental implementation details here; each one determined the sign or magnitude of a reported effect. Studies that combine dataset-native fields with externally imputed or heterogeneous-coverage indicators, or that use rank-based transformations on zero-inflated data, should check both population coverage and clustering stability explicitly, rather than assume a pipeline behaves the same way once population or preprocessing choices change even slightly.

We ran the sensitivity analyses flagged in Section~\ref{sec:method} against alternative $k$, repeated random initializations, retained outliers, exclusion of the AIDev-derived indicators, and the complete-case restriction (Section~\ref{sec:threats}); their results were mixed rather than uniformly reassuring. The clustering is highly stable to different random seeds (Adjusted Rand Index $0.98$--$0.99$) and reasonably stable to the choice of $k$, but considerably less stable to whether outliers are removed (Adjusted Rand Index $0.64$, concentrated in the Senior profile) and least stable of all to whether the three agent-usage indicators are included (Adjusted Rand Index $0.32$), which also weakens RQ1's effect specifically. This last result should temper how RQ1's large effect size is read: part of it plausibly reflects that the clustering was built, in part, from the same kind of agent-usage signal RQ1 measures as an outcome. RQ1 and RQ2 both proved robust to the complete-case restriction -- their group means and significance pattern were nearly unchanged when the 1,755 imputed developers were excluded. Note that we only reran RQ1/RQ2's actual outcome statistics under two of the four checks (the complete-case restriction and the indicator exclusion); the alternative-$k$ and retained-outliers checks assessed the stability of the clustering itself, not whether RQ1/RQ2's conclusions hold under them. RQ1's direction is consistent across every check performed, but its magnitude is the most sensitive to a specific design choice (the indicator exclusion above); a third candidate outcome considered during development but ultimately excluded from this study did not survive the complete-case check at all (Section~\ref{sec:threats}), which is itself a reminder that not every association that looks reportable early in an analysis stays that way once the pipeline is stress-tested.

Bootstrap confidence intervals around every reported effect size (Section~\ref{sec:results}) are consistent with the picture above: RQ1's and RQ2's pairwise intervals are tight relative to their point estimates, while RQ3's interval is comparatively wider, reflecting its much smaller and unevenly sized human-authored comparison group. Future work should extend RQ3 as agent-usage data accumulates for a larger and more representative set of human-authored baseline repositories, rather than only the currently available high-star subset, which would let the human-baseline comparison be assessed with more balanced group sizes than the present analysis permits. More broadly, as coding agents and the norms around reviewing their output continue to change, these associations should be treated as a snapshot of one adoption period rather than a stable characterization of human--agent collaboration on GitHub.
